% TODO make hw stylesheet
\documentclass[10pt]{article}
\usepackage[letterpaper,top=1in,left=1in,right=1in]{geometry}
\usepackage[dvipsnames,svgnames]{xcolor}
\usepackage[shortlabels]{enumitem}
\usepackage[framemethod=TikZ]{mdframed}
\usepackage{amsmath,amssymb,amsthm}
\usepackage[colorlinks]{hyperref}
\usepackage{multicol}
\usepackage{tabularx}
\usepackage{bbm}

\hypersetup{
    colorlinks=true,
    linkcolor=blue,
    filecolor=magenta,
    urlcolor=magenta,
    pdftitle={MAT 549 HW}
}

\usepackage{mathtools}
\usepackage{thmtools}
\usepackage[textsize=footnotesize]{todonotes}
\usepackage{titling}
\usepackage{cancel}

\reversemarginpar
\mdfdefinestyle{mdbluebox}{roundcorner=10pt,innerbottommargin=9pt,
linecolor=blue,backgroundcolor=TealBlue!5,}
\declaretheoremstyle[headfont=\sffamily\bfseries\color{MidnightBlue},
mdframed={style=mdbluebox},]{thmbluebox}

\mdfdefinestyle{mdbloobox}{frametitlefont=\bfseries,innerbottommargin=8pt,
nobreak=true,backgroundcolor=TealBlue!5,linecolor=blue,}
\declaretheoremstyle[headfont=\bfseries\color{MidnightBlue},
mdframed={style=mdbloobox},headpunct={\\[3pt]},postheadspace=0pt,]{thmbloobox}

\mdfdefinestyle{mdredbox}{frametitlefont=\bfseries,innerbottommargin=8pt,
nobreak=true,backgroundcolor=Salmon!5,linecolor=RawSienna,}
\declaretheoremstyle[headfont=\bfseries\color{RawSienna},
mdframed={style=mdredbox},headpunct={\\[3pt]},postheadspace=0pt,]{thmredbox}

\mdfdefinestyle{mdgreenbox}{linecolor=ForestGreen,backgroundcolor=ForestGreen!5,
linewidth=2pt,rightline=false,leftline=true,topline=false,bottomline=false,}
\declaretheoremstyle[headfont=\bfseries\sffamily\color{ForestGreen!70!black},
mdframed={style=mdgreenbox},headpunct={ --- },]{thmgreenbox}

\mdfdefinestyle{mdorangebox}{linecolor=RedOrange,backgroundcolor=RedOrange!5,
linewidth=2pt,rightline=false,leftline=true,topline=false,bottomline=false,}
\declaretheoremstyle[headfont=\bfseries\sffamily\color{RedOrange!70!black},
mdframed={style=mdorangebox},headpunct={ --- },]{thmorangebox}

\mdfdefinestyle{mdblackbox}{linecolor=black,backgroundcolor=RedViolet!5!gray!5,
linewidth=3pt,rightline=false,leftline=true,topline=false,bottomline=false,}
\declaretheoremstyle[mdframed={style=mdblackbox}]{thmblackbox}

\declaretheoremstyle[spaceabove=6pt,spacebelow=6pt]{boldhead}
\declaretheoremstyle[spaceabove=6pt,spacebelow=6pt,headfont=\normalfont\itshape]{italichead}

\declaretheorem[style=boldhead,name=Problem]{problem}
\declaretheorem[style=boldhead,name=Problem,numbered=no]{problem*}
\declaretheorem[style=boldhead,name=Setup]{setup} % for config geo
\declaretheorem[style=boldhead,name=Example,sibling=problem]{example}
\declaretheorem[style=boldhead,name=$\bigstar$ Problem,sibling=problem]{niceprob}
\declaretheorem[style=boldhead,name=Theorem,sibling=problem]{theorem}
\declaretheorem[style=boldhead,name=Corollary,sibling=theorem]{corollary}
\declaretheorem[style=boldhead,name=Proposition,sibling=theorem]{prop}
\declaretheorem[style=boldhead,name=Lemma,numberwithin=problem]{lemma}
\declaretheorem[style=boldhead,name=Theorem,numbered=no]{theorem*}
\declaretheorem[style=boldhead,name=Lemma,numbered=no]{lemma*}
\declaretheorem[style=boldhead,name=Claim,numberwithin=problem]{claim}
\declaretheorem[style=boldhead,name=Claim,numbered=no]{claim*}
\declaretheorem[style=boldhead,name=Remark,numberwithin=problem]{remark}
\declaretheorem[style=boldhead,name=Remark,numbered=no]{remark*}
\declaretheorem[style=boldhead,name=Definition,sibling=theorem]{definition}
\declaretheorem[style=boldhead,name=Definition,numbered=no]{definition*}

\providecommand{\ol}{\overline}
\providecommand{\ceil}[1]{\left\lceil #1 \right\rceil}
\providecommand{\floor}[1]{\left\lfloor #1 \right\rfloor}
\newcommand{\dd}{\,\mathrm{d}}
\providecommand{\ul}{\underline}
\providecommand{\eps}{\varepsilon}
\providecommand{\half}{\frac{1}{2}}
\providecommand{\CC}{\mathbb C}
\providecommand{\FF}{\mathbb F}
\providecommand{\II}{\mathbb I}
\providecommand{\NN}{\mathbb N}
\providecommand{\QQ}{\mathbb Q}
\providecommand{\RR}{\mathbb R}
\providecommand{\ZZ}{\mathbb Z}
\providecommand{\ind}{\mathbbm 1}
\providecommand{\dg}{^\circ}
\providecommand{\ii}{\item}
\providecommand{\setdiff}{\smallsetminus}
\providecommand{\alert}[1]{{\sffamily\textbf{\textcolor{blue}{#1}}}}
\providecommand{\inv}{^{-1}}
\providecommand{\mb}[1]{\mathbf{#1}}
\newcommand{\what}{\widehat}

\newenvironment{soln}{\begin{proof}[Solution]}{\end{proof}}

\providecommand{\pow}{\operatorname{Pow}}
\providecommand{\vocab}[1]{{\textbf{\textcolor{ForestGreen}{#1}}}}
\providecommand{\lcm}{\operatorname{lcm}}
\providecommand{\abs}[1]{\left\lvert #1\right\rvert}
\providecommand{\smabs}[1]{\lvert #1\rvert}
\providecommand{\innerprod}[1]{\langle\!\langle #1\rangle\!\rangle}
\providecommand{\norm}[1]{\left\| #1\right\|}
\providecommand{\smnorm}[1]{\| #1\|}
\providecommand{\gr}{\operatorname{gr}}

\usepackage{mathrsfs}
\providecommand{\tanvec}[1]{\frac{\partial}{\partial #1}}

% place title at top right
\setlength{\droptitle}{-8ex}
\pretitle{\begin{flushright}\Large\bfseries}
\posttitle{\par\end{flushright}}
\preauthor{\begin{flushright}}
\postauthor{\end{flushright}}
\predate{\begin{flushright}}
\postdate{\end{flushright}}

\title{MAT 549 HW02}
\author{\large Aditya Pahuja}
\date{\large Due 02/17}
\begin{document}
\maketitle

\begin{problem}
    Compute the flow of each of the following vector fields on $\RR^2$:
    \begin{enumerate}[(a)]
        \ii $V = y\tanvec x + \tanvec y$
	\ii $W = x\tanvec x + 2y\tanvec y$
	\ii $X = x\tanvec x - y\tanvec y$
	\ii $Y = y\tanvec x + x\tanvec y$
    \end{enumerate}
\end{problem}

\begin{soln}
    (a) Let $\gamma_{(a,b)}(t) = (f(t),g(t))$ be the integral curve
    along $V$ based at the point $(a,b)\in\RR^2$,
    i.e. $\gamma_{(a,b)}(0) = (a,b)$.
    Then (dropping the subscript for brevity), $\gamma$ satisfies
    \[
	\gamma'(t) = f'(t)\tanvec x + g'(t)\tanvec y = V(f(t),g(t))
	= g(t)\tanvec x + \tanvec y.
    \]
    This differential equation has a solution over all of $\RR$, namely
    $g(t) = t + b$ and $f(t) = \half t^2 + bt + a$.
    The resulting flow is
    \begin{equation*}
	\varphi_t^V(a,b) = \gamma_{(a,b)}(t) = \left(\half t^2+bt+a,t+b\right).
    \end{equation*}

    (b) Let $\gamma_{(a,b)}(t) = (f(t),g(t))$ be the integral curve
    along $W$ based at the point $(a,b)\in\RR^2$.
    Then, $\gamma$ satisfies
    \[
	\gamma'(t) = f'(t)\tanvec x + g'(t)\tanvec y = W(f(t),g(t))
	= f(t)\tanvec x + 2g(t)\tanvec y.
    \]
    We can solve the resulting differential equation to get
    $f(t) = ae^t$, $g(t) = be^{2t}$.
    The resulting flow is
    \begin{equation*}
	\varphi_t^W(a,b) = \gamma_{(a,b)}(t) = (ae^t,be^{2t}).
    \end{equation*}

    (c) Let $\gamma_{(a,b)}(t) = (f(t),g(t))$ be the integral curve
    along $X$ based at the point $(a,b)\in\RR^2$.
    Then, $\gamma$ satisfies
    \[
	\gamma'(t) = f'(t)\tanvec x + g'(t)\tanvec y = X(f(t),g(t))
	= f(t)\tanvec x - g(t)\tanvec y.
    \]
    We can solve the resulting differential equation to get
    $f(t) = ae^t$, $g(t) = be^{-t}$.
    The resulting flow is
    \begin{equation*}
	\varphi_t^X(a,b) = \gamma_{(a,b)}(t) = (ae^t,be^{-t}).
    \end{equation*}

    (d) Let $\gamma_{(a,b)}(t) = (f(t),g(t))$ be the integral curve
    along $Y$ based at the point $(a,b)\in\RR^2$.
    Then, $\gamma$ satisfies
    \[
	\gamma'(t) = f'(t)\tanvec x + g'(t)\tanvec y = Y(f(t),g(t))
	= g(t)\tanvec x + f(t)\tanvec y.
    \]
    We can solve the resulting differential equation to get
    $f(t) = \frac{a+b}2\cdot e^x + \frac{a-b}2\cdot e^{-x}$ and
    $g(t) = \frac{a+b}2\cdot e^x - \frac{a-b}2\cdot e^{-x}$.
    The resulting flow is
    \begin{equation*}
	\varphi_t^Y(a,b) = \gamma_{(a,b)}(t)
	= \left(\frac{a+b}2\cdot e^x + \frac{a-b}2\cdot e^{-x},
	\frac{a+b}2\cdot e^x - \frac{a-b}2\cdot e^{-x}\right).
    \end{equation*}
\end{soln}
\pagebreak

\begin{problem}
    For any integer $n\ge 1$, define a flow on
    the odd-dimensional sphere $S^{2n-1}\subseteq\CC^n$ by
    $\theta(t,z) = e^{it}z$. Show that the infinitesimal generator
    of $\theta$ is a smooth nonvanishing vector field on $S^{2n-1}$.
\end{problem}

\begin{soln}
    The infinitesimal generator of $\theta(t,z)$ is the vector field
    $V(z) = {\theta^{(z)}}'(0) = iz$.
    Since no point on $S^{2n-1}$ is zero,
    $iz$ is also never zero, so this is indeed
    a nonvanishing smooth vector field on $S^{2n-1}$.
\end{soln}

\begin{problem}
    Suppose $M$ is a smooth, compact manifold that admits
    a nowhere vanishing vector field. Show that there exists
    a smooth map $F\colon M\to M$ that is homotopic to the identity
    and has no fixed points.
\end{problem}

\begin{soln}
    We will show that for some $t>0$
    the flow $\varphi_t^V\colon M\to M$ of $M$ along $V$
    satisfies the desired properties.
    Since the flow $\varphi_0^V$ at $t=0$ is the identity map,
    the flow at time $t$ is homotopic to the identity for every $t$,
    so we only need to find a $t$ for which
    $\varphi_t^V$ has no fixed points.

    Let $p\in M$ be an arbitrary point
    and $(x_1,\dots,x_m)$ a system of local coordinates
    on a neighborhood $U$ of $p$ such that
    the coordinates of $p$ are $(0,\dots,0)$ and $V(p) = \tanvec{x_1}$.
    We will alter the given coordinates to force
    $V\equiv\tanvec{x_1}$ in a neighborhood of $p$.
    Let $\Sigma$ be the hypersurface given by the points inside $U$
    with $x_1 = 0$; clearly $\Sigma$ is transverse to $V$ at $p$,
    since $\tanvec{x_1}$ is not tangent to $\Sigma$.

    Now, define $\Phi\colon(-\eps,\eps)\times\Sigma\to M$
    by $\Phi(t,q) = \varphi_t^V(q)$, i.e.\ $\Phi$ is
    the flow of $\Sigma$ along $V$;
    $\dd\Phi_{(0,p)}$ is invertible because
    its derivative with respect to $t$ at that point is just $V(p)$
    and the differential with respect to $q$ is $T_p\Sigma$
    which, together with $V(p)$, spans $T_pM$.
    The upshot of this is that some neighborhood of $(0,p)$
    gets mapped diffeomorphically to a neighborhood of $p$, so
    take a system of local coordinates $(y_1,\dots,y_m)$
    on the former neighborhood with $y_1 = t$
    and $(y_2,\dots,y_m)=q$ coordinates for $\Sigma$ at $p$.
    Via $\Phi$, these map to local coordinates $(\tilde y_1,\dots,\tilde y_m)$
    on a neighborhood of $p$ in which $V\equiv\tanvec{x_1}$,
    since flowing in the direction of $V$ corresponds to
    moving in the $y_1$ direction.
    
    What we have done is construct local coordinates in
    a neighborhood $U_p$ of $p$
    for which flowing along $V$ corresponds to translation,
    so there exists $t_p > 0$ such that
    $\varphi_t^V$ has no fixed points in $U_p$ for any $0 < t < t_p$.
    Since $M$ is compact, there exist finitely many points
    $p_1,\dots,p_k\in M$ such that $U_{p_1}\cup\cdots\cup U_{p_k} = M$,
    hence $\varphi_{t_0}^V$ has no fixed points
    where $t_0$ is chosen such that $0 < t_0 < \min_{1\le i\le k}t_{p_i}$.
\end{soln}











\end{document}
